\documentclass[11pt]{article}

\usepackage[margin=1in]{geometry}
\usepackage{parskip}
\usepackage{hyperref}
\usepackage{xcolor}
\usepackage{listings}
\usepackage{booktabs}
\usepackage{longtable}
\usepackage{enumitem}
\usepackage{amsmath}

\hypersetup{colorlinks=true, linkcolor=blue, urlcolor=blue, citecolor=blue}

\definecolor{codebg}{rgb}{0.96,0.96,0.96}
\lstset{
  basicstyle=\ttfamily\small,
  backgroundcolor=\color{codebg},
  breaklines=true,
  columns=fullflexible,
  keywordstyle=\color{blue!70!black}\bfseries,
  commentstyle=\color{green!40!black},
  stringstyle=\color{purple},
  frame=single,
  framesep=3pt,
  xleftmargin=3pt,
  xrightmargin=3pt,
  tabsize=2,
}
\lstdefinelanguage{PyShort}{language=Python}

\title{How AutoNixNan Works: A Code-Level Account}
\author{ganesh@cs.utah.edu \\ \small (written with Claude Code)}
\date{\today}

\begin{document}
\maketitle

\begin{abstract}
This document explains, from the source code, how \texttt{bin/autonixnan}
(``AutoNixNan'') and the \texttt{nixnan.so} NVBit tool it drives actually
work, and how its two newer siblings, \texttt{iter-autonixnan} and
\texttt{ksel-autonixnan} (Part III), extend it with different
kernel-scoping strategies. It is not a user guide; it is a line-level
account of the mechanism, written after reading every file involved
(\texttt{bin/autonixnan}, \texttt{src/nixnan.cu}, and the eleven
supporting headers/translation units under \texttt{src/}), and after
empirically exercising the tools against several real CUDA programs --
from small, purpose-built kernels through a full HeCBench port, a
PyTorch/cuSOLVER routine, a toy GPT, a recurrent network, and (Part III,
Section~\ref{sec:limitations}) a 7-billion-parameter, 4-bit-quantized
production LLM. Where the code's behavior diverges from what its own
comments or the project README claim, that divergence is called out
explicitly and is backed by a repeatable observation, not speculation.
\end{abstract}

\tableofcontents

\section{The Big Picture}

AutoNixNan is not one program; it is two, layered:

\begin{enumerate}
\item \textbf{\texttt{nixnan.so}} (source: \texttt{src/*.cu}, \texttt{src/*.cuh},
  \texttt{src/*.h}) --- an NVBit tool. NVBit (NVIDIA Binary Instrumentation
  Tool) is preloaded into a CUDA process via \texttt{LD\_PRELOAD} and gets
  a callback on \emph{every} kernel launch, before that kernel's SASS
  (compiled machine code) ever executes on the GPU. But the SASS rewrite
  itself is not per-launch: \texttt{instrument\_function} (\texttt{src/nixnan.cu})
  guards it with an \texttt{std::unordered\_set<CUfunction>} of functions
  already seen, so it inserts calls to its own device-side functions after
  (and sometimes before) every floating-point instruction only the
  \emph{first} time a given kernel function is launched. Every later
  launch of that same kernel hits the callback again, but
  \texttt{instrument\_function} short-circuits immediately and NVBit just
  runs the already-instrumented SASS via \texttt{nvbit\_enable\_instrumented}
  --- no re-rewriting occurs. This is how nixnan.so inspects the actual
  bits flowing through the GPU's registers at run time without paying the
  instrumentation cost on every single launch.
\item \textbf{\texttt{bin/autonixnan}} (a 420-line Python 3 script) --- a
  wrapper that never touches SASS or the GPU directly. It only launches
  the target program multiple times, each time with a different
  combination of environment variables that steer \texttt{nixnan.so},
  and stitches the resulting log files into one report.
\end{enumerate}

Everything \texttt{autonixnan} does, it does by shelling out to the
target program with \texttt{LD\_PRELOAD=nixnan.so} and a specific
environment. There is no IPC between the Python process and the GPU
process beyond environment variables in and a text log file out. This
matters: every quirk of \texttt{autonixnan}'s behavior traces back to
either (a) what environment variable combination it chose to set for a
given phase, or (b) what \texttt{nixnan.so} actually does --- which is
sometimes \emph{not} what its own \texttt{--help} banner or
\texttt{autonixnan}'s comments say it does (Section~\ref{sec:limitations}) ---
and can also change between checkouts, as Section~\ref{sec:timeout-reenabled}
found out the hard way.

\section{Part I --- \texttt{nixnan.so}}

\subsection{Initialization: \texttt{nvbit\_at\_init}}
\label{sec:init}

NVBit calls \texttt{nvbit\_at\_init()} once, before the target's
\texttt{main()} runs, but after CUDA driver initialization has begun
enough that reading environment variables is possible. \texttt{nixnan.cu}
uses this hook to read every environment variable it supports via the
NVBit-provided \texttt{GET\_VAR\_INT}/\texttt{GET\_VAR\_STR} macros
(defined in \texttt{nvbit\_release/core/nvbit.h}). Each call to one of
these macros \emph{also prints a banner line} to the log stream --- this
is the ``NVBit core environment variables'' banner every instrumented run
of this session emitted. That banner is not decorative: it is a direct
enumeration of every env var the running binary actually reads, in the
order it reads them. \textbf{If a variable does not appear in the banner,
the running binary does not consult it, full stop} --- this is the
diagnostic Section~\ref{sec:timeout-reenabled} uses to confirm
\texttt{NIXNAN\_TIMEOUT} is (now) genuinely read, not just to find dead
code.

Table~\ref{tab:envvars} lists every variable read by \texttt{nixnan.cu}
itself, in banner order, plus two read by a header it includes
(\texttt{fp-histogram.cu}) during \texttt{nvbit\_at\_init}'s call to
\texttt{nixnan::fp\_histogram::init()}.

\begin{longtable}{p{3.2cm}p{2.3cm}p{8.2cm}}
\toprule
\textbf{Variable} & \textbf{Default} & \textbf{Effect} \\
\midrule
\endhead
\texttt{INSTR\_BEGIN} / \texttt{INSTR\_END} & 0 / UINT32\_MAX & Instruction-index interval to instrument (unused by autonixnan). \\
\texttt{TOOL\_VERBOSE} & 0 & Extra diagnostic printing during instrumentation setup. \\
\texttt{ENABLE\_FUN\_DETAIL} & 0 & Print the full mangled kernel name every time a kernel re-runs, not just the first time. \\
\texttt{PRINT\_ILL\_INSTR} & 0 & (declared, not read further in the code paths exercised here). \\
\texttt{SAMPLING} & 0 & Use the already-instrumented SASS only on every $N$th launch of a repeated kernel (others run the original, uninstrumented version); Section~\ref{sec:sampling-launch-gate}. \\
\texttt{INSTRUMENT\_EXCEPTIONS} & 1 & Master switch for the NaN/Inf/Subnormal/Div-by-0 instrumentation described in Section~\ref{sec:exceptions}. \\
\texttt{INSTR\_MEM} & 0 & Also instrument 32/64/128-bit global-memory stores (\texttt{STG}) for NaN, Section~\ref{sec:memcheck}. \\
\texttt{HISTOGRAM} & 0 & Set implicitly to 1 if \texttt{BIN\_SPEC\_FILE} is non-empty; enables the exponent-histogram/bin-threshold mechanism, Section~\ref{sec:histogram}. \\
\texttt{BIN\_SPEC\_FILE} & (empty) & Path to a JSON file describing which exponent ranges, per format, should be watched; auto-created as a template and the process exits if the file does not yet exist. \\
\texttt{NIXNAN\_TIMEOUT} & 0 (disabled) & Seconds before the process gracefully self-terminates (\texttt{nnterminate}); confirmed active, Section~\ref{sec:timeout-reenabled}. \\
\texttt{LOGFILE} & (empty) & Redirect all \texttt{nnout()} output (everything prefixed \texttt{\#nixnan:}) from \texttt{stderr} to this file. \\
\texttt{LINE\_INFO} & 1 & Attempt to resolve source file:line for each instrumented instruction (wrapped in a \texttt{SIGSEGV} handler --- see \texttt{recording.cu}). \\
\texttt{LOG\_KERNELS} & (empty) & Path to a log file; if set, \textbf{disables all instrumentation} and instead just logs every kernel launch and its wall-clock execution time. Section~\ref{sec:logkernels}. \\
\texttt{MAX\_ERRORS} & 0 (unlimited) & Self-terminate (\texttt{SIGINT} via \texttt{nnterminate}) after this many exception reports. \\
\texttt{TIME\_KERNELS} & 0 & Print a ``Kernel [\ldots] execution time: N microseconds'' line for every kernel launch, independent of \texttt{LOG\_KERNELS}. \\
\texttt{FUNCTION\_WHITELIST} & (empty) & Path to a newline-separated file of function names; only these are instrumented. \\
\texttt{FUNCTION\_BLACKLIST} & (empty) & Complement of the above; mutually exclusive with it (the process exits with an error if both are set). \\
\bottomrule
\caption{Environment variables consulted by the running \texttt{nixnan.so} binary, in the order \texttt{nvbit\_at\_init} reads them (as of the current \texttt{print\_histo\_instrn} checkout). Absence from this table means the variable is \emph{not currently wired up} -- confirmed by grep against the current source, not just by a variable's absence from a possibly-stale banner capture.}
\label{tab:envvars}
\end{longtable}

\subsection{Deciding what to instrument}

For every kernel launch, NVBit calls
\texttt{nvbit\_at\_cuda\_event(ctx, is\_exit, cbid, \ldots)}, once on entry
and once on exit, for the CUDA driver calls that launch work
(\texttt{cuLaunchKernel} and its cooperative variants). On entry
(\texttt{is\_exit == false}):

\textbf{Update:} the code quoted in this subsection for the launch
callback and \texttt{should\_instrument} was accurate for an earlier
checkout, but \texttt{SAMPLING} has since been refactored onto a
different, and simpler, mechanism --- verified directly against the
current \texttt{print\_histo\_instrn} source, not assumed from the
earlier text. The corrected code and explanation follow; see also
Section~\ref{sec:sampling-launch-gate} for what this means for
\texttt{SAMPLING} specifically.

\begin{lstlisting}[language=C++]
std::string kernel_name = nvbit_get_func_name(ctx, p->f);
...
bool enable_instr = instrument_function(ctx, p->f);
if (analyzed_kernels[kernel_name] == 0) {
  nnout() << "running kernel [" << kernel_name << "] ...";
} else if (func_details) {
  nnout() << "running kernel [" << kernel_name << "] ...";
}
auto kernel_count = analyzed_kernels[kernel_name]++;
nvbit_enable_instrumented(ctx, p->f,
    sampling > 0 ? kernel_count % sampling == 0 : true);
\end{lstlisting}

Two things changed from the earlier reading of this code. First, the
``running kernel'' print is no longer gated on \texttt{enable\_instr} at
all --- it fires on a kernel's first-ever launch unconditionally, and
also on every later launch if \texttt{ENABLE\_FUN\_DETAIL} (the
\texttt{func\_details} flag) is set, whether or not that particular
launch ends up instrumented. Second, and more importantly, the final
argument to \texttt{nvbit\_enable\_instrumented} is no longer a bare
\texttt{true} --- this is where \texttt{SAMPLING} now lives, discussed
below.

\texttt{instrument\_function} walks the kernel plus every ``related
function'' NVBit knows it might call (\texttt{nvbit\_get\_related\_functions}
--- this includes every template instantiation the compiler linked in as a
\emph{possible} callee, which is why a single cuBLAS \texttt{nrm2\_kernel}
launch can enumerate dozens of unrelated-looking function names; see
Section~\ref{sec:logkernels}). For each such function it calls
\texttt{should\_instrument}, which is where whitelist/blacklist filtering
and --- critically --- \texttt{LOG\_KERNELS} mode converge. On the current
checkout, \texttt{should\_instrument} carries \emph{no} sampling logic at
all:

\begin{lstlisting}[language=C++]
bool should_instrument(CUcontext ctx, CUfunction f) {
  std::string func_name = nvbit_get_func_name(ctx, f);
  bool enable_instr = false;
  if (!kernel_whitelist.empty()) {
    enable_instr = kernel_whitelist.count(func_name);
  } else if (!kernel_blacklist.empty()) {
    enable_instr = !kernel_blacklist.count(func_name);
  } else {
    enable_instr = true;
  }
  enable_instr &= !kernel_logging_enabled;   // <-- LOG_KERNELS forces this to false
  return enable_instr;
}
\end{lstlisting}

So \texttt{LOG\_KERNELS} does not merely add logging on top of the normal
instrumentation --- it hard-disables \texttt{should\_instrument} for
\emph{every} function, meaning \texttt{instrument\_function} inserts zero
device-side calls anywhere. The kernel still runs, at essentially native
speed, because \texttt{nvbit\_\allowbreak enable\_\allowbreak instrumented}
is still called on
every launch (per the callback above) and, with nothing ever instrumented
in the first place, simply runs the (uninstrumented, in this case
identical-to-original) version regardless of what its \texttt{SAMPLING}
argument evaluates to.

\subsubsection*{Where \texttt{SAMPLING} actually lives now}
\label{sec:sampling-launch-gate}

\texttt{SAMPLING} is checked once per launch, in the callback above, as
the boolean argument to
\texttt{nvbit\_\allowbreak enable\_\allowbreak instrumented}:
\begin{lstlisting}[language=C++]
sampling > 0 ? kernel_count % sampling == 0 : true
\end{lstlisting}
This is a structurally different mechanism from
\texttt{should\_\allowbreak instrument}'s whitelist, blacklist, and
\texttt{LOG\_KERNELS} checks, and the two are decoupled:
\texttt{should\_\allowbreak instrument} (via
\texttt{instrument\_\allowbreak function}) decides, once per kernel
\emph{function}, the first time it is ever seen, whether its SASS gets
rewritten with device-side checks inserted at all; \texttt{SAMPLING}, by
contrast, is re-evaluated on \emph{every launch} of that function and
decides whether \emph{this specific launch} runs the already-rewritten
instrumented version or the original uninstrumented one --- a launch-time
toggle between two SASS variants, not a rewrite decision.

This also means the key-mismatch bug \texttt{lu\_solve\_gpu\_report.tex}
documented (the sampling lookup keyed on a function's full mangled
signature while the counter it read was incremented under a
\texttt{cut\_kernel\_name}-truncated name, so the two never matched) does
not exist in this code path: \texttt{kernel\_count} and the modulo check
both key off the same \texttt{kernel\_name} variable, read once at the top
of the callback --- there is no truncated name anywhere in this
mechanism. For a kernel launched $L$ times (\texttt{kernel\_count} running
$0,\dots,L{-}1$), \texttt{SAMPLING=}$N$ enables instrumentation exactly
where \texttt{kernel\_count \% }$N$\texttt{ == 0} --- e.g.\ against
\texttt{lu\_solve\_gpu}'s $L{=}19$-launch loop kernels, \texttt{SAMPLING=16}
should enable only \texttt{kernel\_count}~0 and~16, two of the nineteen.

An open question, not resolved here: running \texttt{lu\_solve\_\allowbreak
gpu} with \texttt{SAMPLING=16} against
\texttt{bin\_spec\_\allowbreak whole.json} reduced total
threshold-report volume by only $\approx\!20\%$ (131$\to$105 reports for
\texttt{matrix.csv}, 111$\to$97 for \texttt{bad\_matrix.csv}) --- much
smaller than the $\approx\!85$--$90\%$ reduction the two-enabled-launches-
out-of-nineteen picture above would suggest if report volume scaled
roughly with each launch's share of the total work (LU elimination's
per-iteration cost shrinks as the remaining submatrix shrinks, so the
$\mathtt{kernel\_count}{=}0$ launch alone touches the full
$20\times20$ matrix while later launches touch a shrinking corner of it
--- plausibly enough to explain \emph{some} of the gap, but this session
did not verify it directly, e.g.\ by instrumenting in a way that
attributes individual reports back to a specific launch index).

For every instruction NVBit hands back for a function that \emph{is}
selected, \texttt{instruction\_info::get\_reginfo(instr)} decides whether
the instruction is floating-point-interesting at all, and if so, what its
operands look like.

\subsection{The opcode-to-operand table}

\texttt{get\_reginfo} (in \texttt{instruction\_info.cu}) matches the
instruction's opcode string against a table of regex patterns baked into
the binary as \texttt{src/instructions.h} --- a byte array
(\texttt{instructions\_json}, 5{,}433 bytes) that is really just a JSON
document, embedded at compile time and parsed once at first use with
\texttt{nlohmann::json}. Two representative entries:

\begin{lstlisting}
"HMMA.16816.F16": {
  "registers": [
    {"type": "f16", "count": 4},
    {"type": "f16", "count": 8},
    {"type": "f16", "count": 4},
    {"type": "f16", "count": 4}
  ],
  "description": "Half-precision matrix multiply accumulate into half-precision accumulator"
},
"DFMA(\\.RZ|\\.RP|\\.RM)?(\\.FTZ)?(?:\\.(RP|RM|RZ))?(?:\\.SAT)?": {
  "registers": [
    {"type": "f64", "count": 1},
    {"type": "f64", "count": 1},
    {"type": "f64", "count": 1}
  ],
  "description": ""
}
\end{lstlisting}

Each entry says, per operand position, which floating-point format the
operand is (\texttt{f16}/\texttt{bf16}/\texttt{f32}/\texttt{f64}) and how
many scalar values are packed into it (e.g.\ a tensor-core \texttt{HMMA}
operand carrying 8 packed half-precision values). For each such operand,
\texttt{get\_regs} figures out \emph{how} to actually read the value at
instrumentation-insertion time --- a register, a uniform register, or a
constant-bank (\texttt{cbank}) immediate --- and builds a list of
zero-argument closures (\texttt{reginsertion}, i.e.\
\texttt{std::function<void()>}) that, when called later, will emit the
NVBit call
(\texttt{nvbit\_add\_call\_arg\_reg\_val}/\texttt{\_ureg\_val}/\texttt{\_cbank\_val})
that actually smuggles that operand's bits into the injected device
function call. A literal double/float immediate operand
(\texttt{IMM\_DOUBLE}) is instead classified right there at instrumentation
time via host-side \texttt{fpclassify} and, if already NaN/Inf/subnormal,
produces an immediate warning line rather than a runtime check (there is
nothing to check at run time --- the value never changes).

This table is also where the 24-line \texttt{is\_f64high} check matters:
some SASS forms split a 64-bit double across two 32-bit registers and
give the \emph{high} half its own opcode suffix (\texttt{...64H}); the
table's register-start-offset logic (\texttt{reg\_start = op->u.reg.num -
(f64high ? 1 : 0)}) exists purely to correctly pair the two halves back
into one 64-bit value.

\subsection{Exception detection, end to end}
\label{sec:exceptions}

For every FP-classified instruction (and \texttt{INSTRUMENT\_EXCEPTIONS} is
1 by default), \texttt{instrument\_function} in \texttt{nixnan.cu} inserts
\emph{two} calls to the same device function, \texttt{nixnan\_check\_regs}
(defined in \texttt{inject\_funs.cu}): once \texttt{IPOINT\_AFTER} (reading
the operand \emph{after} it takes its post-instruction guard-predicate
value --- effectively the output/destination operands) and once
\texttt{IPOINT\_BEFORE} (the input operands). Each call carries, as
variadic arguments, a packed \texttt{reginfo} struct (a bit-packed 32-bit
struct --- \texttt{count:6, type:4, half\_h0:1, half\_h1:1, div0:1,
operand:2, num\_regs:4, f64high:1}, statically asserted to fit in one
\texttt{int32\_t}) followed by the actual register values.

\texttt{nixnan\_check\_regs} then, per operand, classifies the raw bits
using the IEEE-754 bit-pattern predicates in \texttt{fp\_utils.cuh}
(\texttt{half\_is\_nan}, \texttt{float\_is\_inf}, \texttt{double\_is\_subnorm},
etc. --- pure exponent/mantissa bit masking, no floating-point arithmetic
at all, so this cannot itself introduce rounding artifacts):

\begin{lstlisting}[language=C++]
uint32_t float_is_nan(uint32_t f) {
  return (float_exp(f) == 0xFF && float_mant(f) != 0) ? E_NAN : 0;
}
uint32_t float_is_inf(uint32_t f) {
  return (float_exp(f) == 0xFF && float_mant(f) == 0) ? sgninf(f) : 0;
}
\end{lstlisting}

The five exception bits are \texttt{E\_NAN}, \texttt{E\_INF},
\texttt{E\_NINF} (negative infinity, reported separately from positive),
\texttt{E\_SUB} (subnormal), and \texttt{E\_DIV0} --- the last of which is
really just ``this operand is exactly zero'', only checked for operands
the instruction table flags with \texttt{"div0": true} (i.e.\ divisor
operands of division-shaped instructions). \texttt{half2}/\texttt{bf162}
variants unpack two packed half-precision lanes from one 32-bit register
before classifying each.

The per-lane exception bitmask is then combined across the whole warp with
a \texttt{\_\_shfl\_sync} reduction (any lane's exception counts for the
warp), and only the warp's first active lane does the recording, to avoid
32-way redundant writes:

\begin{lstlisting}[language=C++]
if (first_laneid == laneid) {
  for (int op = 0; op < OPERANDS; op++) {
    uint32_t exce = exces[op];
    if (exce != 0) {
      uint32_t num_exceptions = recorder.record(inst_id, exce, op);
      if (num_exceptions == 0) {
        // first time this (instruction, operand) pair has ever
        // hit this exact exception -- push a live report
        for (int skip = 0; skip < 2; skip++) {
          exception_info ei(get_ctaid(), get_warpid(), inst_id, exce, op, UNKNOWN, skip);
          channel_dev->push(&ei, sizeof(exception_info));
        }
      }
    }
  }
}
\end{lstlisting}

\texttt{recorder.record} (\texttt{device\_recorder::record},
\texttt{recording.h}) is an \texttt{atomicAdd} into a flat device array,
indexed by \texttt{(instruction id, operand, exception bitmask)}. This is
the counter that \textbf{every} occurrence increments, whether or not a
message is pushed anywhere --- it is the source of the ``repeats'' numbers
in the final report. The channel push, by contrast, only happens the
\emph{first} time a given \texttt{(instruction, operand, exception-bits)}
triple is ever seen (\texttt{num\_exceptions == 0} means this was the very
first increment) --- this is what produces the live
\texttt{\#nixnan: error [\ldots] detected in operand \ldots} lines this
session saw streamed to the console, and it happens exactly once per
distinct site, no matter how many times that site actually fires.

\subsubsection*{The channel and receiver thread}

NVBit tools cannot easily do host-side I/O from device code, so
\texttt{nixnan.so} uses NVBit's lock-free \texttt{ChannelDev}/\texttt{ChannelHost}
ring buffer: device code \texttt{push}es fixed-size \texttt{exception\_info}
records; a dedicated host \texttt{std::thread} (\texttt{recv\_thread\_fun}
in \texttt{nixnan.cu}) polls \texttt{channel\_host.recv} and, for each
record, looks up the recorded instruction's SASS text, function name, and
source location (all previously stashed by \texttt{recorder::mk\_entry} at
instrumentation time) and prints the formatted error line. Every push
happens \emph{twice}, with a ``skip'' flag on the second copy --- this is
an NVBit idiom to guarantee the channel has fully drained even if the
consuming thread is mid-poll (\texttt{to\_skip()} entries are silently
discarded by the receiver).

\subsubsection*{The final aggregate report}

At context teardown (\texttt{nvbit\_at\_ctx\_term}), the device-side error
array is copied back to the host and folded into per-\texttt{(format,
is\_memory)} totals. For each of the five exception kinds, the code
computes \emph{two} numbers from the same per-instruction/operand counter
table:

\begin{lstlisting}[language=C++]
if (exce & error) {
  std::get<0>(exception_counts[{type, is_mem}][index]) += errors;   // total occurrences (repeats, so far)
  if (errors > 0) std::get<1>(exception_counts[{type, is_mem}][index])++;  // one more distinct site
}
...
// later: subtract the per-format unique-site count out of the running total
std::get<0>(exception_counts[{ftype, is_mem}][i]) -= std::get<1>(exception_counts[{ftype, is_mem}][i]);
\end{lstlisting}

That subtraction is exactly why the printed report reads
\texttt{NaN: 22 (154 repeats)}: 22 is the number of distinct
(instruction, operand) sites that ever produced a NaN even once, and 154
is \emph{additional} occurrences beyond each site's first (i.e.\ total
occurrences minus the 22 ``first hits'' already counted once each in the
unique-site tally) --- so \texttt{unique + repeats} is the true total
number of times that exception fired, summed over every site. This is
the exact metric this session used to compare \texttt{rd\_nixnan.cu}
against its 10x-grid-point variant: unique-site count is a static
property of the kernel's source code and stays fixed as the problem
scales, while repeats measures how much bad floating-point behavior an
actual run produced, and that number does scale with problem size.

\subsection{Memory-instruction checking}
\label{sec:memcheck}

If \texttt{INSTR\_MEM=1}, every global-memory store (\texttt{STG}, any
width) additionally gets a call to \texttt{nixnan\_check\_nans}
(\texttt{memcheck.cu}) inserted \texttt{IPOINT\_BEFORE}.
\texttt{meminstrumentation.cu} first has to work out \emph{what type} is
being stored, since a raw \texttt{STG} instruction's operand is just an
untyped register: it walks backwards through the current basic block
looking for the most recent instruction that wrote that register and
reuses \emph{that} instruction's already-known operand type
(\texttt{find\_type}). If no typed producer is found, the store is
skipped rather than guessed at. Unlike the exception path above, the
memory-instruction checker only ever reports \texttt{NaN} (not
Inf/Subnormal/Div0) --- this is why every ``\texttt{\ldots Memory
Operations}'' block in the report only ever prints a \texttt{NaN:} line.

\subsection{The exponent-histogram / bin-threshold mechanism}
\label{sec:histogram}

This is the machinery behind \texttt{HISTOGRAM}/\texttt{BIN\_SPEC\_FILE},
and it is what \texttt{autonixnan}'s ``extreme-value scan'' phase
(Section~\ref{sec:phase2}) is actually built on. It runs as a
\emph{parallel, independent} instrumentation pass alongside the exception
checker above --- every FP instruction that gets a
\texttt{nixnan\_check\_regs} call also, unconditionally, gets a call to
\texttt{nixnan\_fp\_histogram\_counter} inserted (\texttt{fp-histogram.cu});
whether that call does anything is gated at run time by
\texttt{histogram\_enabled} (true if \texttt{HISTOGRAM=1} or
\texttt{BIN\_SPEC\_FILE} is set).

\texttt{BIN\_SPEC\_FILE} is a small JSON document naming, per format, a
list of \texttt{[lower, upper]} \emph{unbiased signed exponent} ranges to
watch, plus a \texttt{count} (report every \texttt{count}-th value landing
in a watched bin) and optional \texttt{max\_reports} (self-terminate a
kernel-function's reporting after this many). If the named file does not
exist, \texttt{nixnan.so} writes an empty template there and
\texttt{exit(0)}s --- by design, the first invocation with a
not-yet-existing spec file is a no-op setup step, not a real run.

At run time, \texttt{record()} (in \texttt{fp-fun.cu}) does two things
for every FP value seen: it always \texttt{atomicAdd}s into a flat
\texttt{[format][exponent]} histogram array (this is what powers the
``Exponent range for f16: [-14, 15]'' summary line printed at
\texttt{term()}, when \texttt{HISTOGRAM=1} is used directly rather than
through \texttt{autonixnan}), and, if that exponent falls in one of the
function's configured bins, increments a per-bin counter and --- every
\texttt{count}-th time --- pushes a \texttt{bin has reached threshold}
report through its own separate channel/receiver-thread pair
(structurally identical to the exception path, but entirely separate
state). \textbf{Crucially, the valid exponent range for each format
explicitly excludes zero and infinity} (see \texttt{FP\_FORMAT\_RANGES} in
\texttt{bin/autonixnan}, and \texttt{bin\_from\_json}'s range validation in
\texttt{fp-histogram.cu}, which rejects any bin touching the
all-zeros/all-ones exponent encodings) --- \textbf{a value that jumps
straight from a normal, finite magnitude to \texttt{Inf} in a single
instruction never lands in any bin at all, and is therefore invisible to
this mechanism.} This is exactly what this session found empirically
with the mixed-precision-Cholesky Phase-D overflow experiment
(Section~\ref{sec:limitations}): the scan flagged an unrelated fp64 QR
kernel while missing the real, program-confirmed FP16 overflow entirely,
because that overflow was a one-step jump to \texttt{Inf} with no
gradually-climbing finite exponent trail leading up to it.

\subsection{\texttt{LOG\_KERNELS} mode}
\label{sec:logkernels}

As established above, \texttt{LOG\_KERNELS} disables all instrumentation.
What it \emph{does} do: at kernel-exit time,
\texttt{nvbit\_at\_cuda\_event} prints
\texttt{Kernel [name] execution time: N microseconds} for every launch
(gated on \texttt{time\_kernels || kernel\_logging\_enabled}, so this
line's presence does not actually require \texttt{TIME\_KERNELS=1} once
\texttt{LOG\_KERNELS} is set --- \texttt{autonixnan}'s
\texttt{collect\_kernels()} relies on exactly this). Separately,
\texttt{instrument\_function} also prints one
\texttt{Function [related] of kernel [top-level]} line for \emph{every}
NVBit-enumerated related function of a kernel, which for a library-heavy
binary (cuBLAS/cuSOLVER template instantiations) can be enormous --- an
n=4096 cuSOLVER Cholesky factorization run in this session produced over
10{,}000 such lines, none of which match \texttt{autonixnan}'s
\texttt{\#nixnan: Kernel [\ldots] execution time:} parser regex and are
therefore silently ignored by \texttt{\_parse\_kernel\_log}. They are
harmless noise from \texttt{autonixnan}'s point of view, but they do
inflate the log file and are worth knowing about if reading a
\texttt{LOG\_KERNELS} log by hand.

\paragraph{An extra bit of detail: no first-launch cost here.} Run
\texttt{lu\_solve\_gpu} again, this time with \emph{only}
\texttt{LOG\_KERNELS} set (to a log file path; \texttt{TIME\_KERNELS} left
at its default, relying on the \texttt{kernel\_logging\_enabled} half of
the gate above) --- same binary, same input, same 8 kernels as
Section~\ref{sec:time-kernels-first-launch}'s table:

\begin{center}
\small
\begin{tabular}{lrrr}
\toprule
Kernel & Launches & First launch & Typical (median) \\
\midrule
\texttt{init\_identity\_kernel} & 2 & 5{,}006~\textmu s & 31~\textmu s \\
\texttt{find\_pivot\_kernel} & 19 & 21~\textmu s & 10~\textmu s \\
\texttt{swap\_rows\_kernel} & 19 & 16~\textmu s & 10~\textmu s \\
\texttt{compute\_multipliers\_kernel} & 19 & 15~\textmu s & 10~\textmu s \\
\texttt{eliminate\_kernel} & 19 & 14~\textmu s & 9~\textmu s \\
\texttt{apply\_permutation\_kernel} & 1 & 25~\textmu s & --- \\
\texttt{forward\_substitution\_kernel} & 1 & 15~\textmu s & --- \\
\texttt{backward\_substitution\_kernel} & 1 & 14~\textmu s & --- \\
\bottomrule
\end{tabular}
\end{center}
No $\approx\!300$ms first-launch spike anywhere in this table --- every
kernel's first launch is within a small constant factor of its own later
launches (the one partial exception, \texttt{init\_identity\_kernel}'s
5{,}006~\textmu s vs.\ its own 31~\textmu s second launch, is ordinary
CUDA-driver/context warm-up on the very first kernel launch of the whole
process, not SASS rewriting --- it is two orders of magnitude smaller than
any first-launch cost in Section~\ref{sec:time-kernels-first-launch}'s
table, and every other kernel here shows no such spike at all). This is
exactly the confirming half of that section's causal claim: under
\texttt{LOG\_KERNELS}, \texttt{should\_instrument} forces
\texttt{enable\_instr} to false for every function
(\texttt{enable\_instr \&= !kernel\_logging\_enabled}), so
\texttt{instrument\_function} never rewrites any kernel's SASS at all ---
and with no rewrite to pay for, the first-launch cost simply does not
appear.

\subsection{\texttt{TIME\_KERNELS} Under Full Instrumentation: the First-Launch Cost}
\label{sec:time-kernels-first-launch}

The previous subsection covers \texttt{TIME\_KERNELS} under
\texttt{LOG\_KERNELS} mode, where instrumentation is disabled entirely.
Run \texttt{TIME\_KERNELS=1} on its own instead --- \texttt{LD\_PRELOAD}
and \texttt{TIME\_KERNELS=1}, nothing else, so
\texttt{INSTRUMENT\_EXCEPTIONS} stays at its default (on) --- against
\texttt{lu\_solve\_gpu} (this repository's \texttt{lu\_solve\_gpu/},
well-conditioned \texttt{matrix.csv}) and every one of its 8 kernels'
timings tells the same story, shown in
Table~\ref{tab:time-kernels-first-launch}:

\begin{table}[!htb]
\centering
\small
\begin{tabular}{lrrrr}
\toprule
Kernel & Launches & First launch & Typical (median) & Ratio \\
\midrule
\texttt{init\_identity\_kernel} & 2 & 312{,}451~\textmu s & 21~\textmu s & $\approx\!14{,}900\times$ \\
\texttt{find\_pivot\_kernel} & 19 & 298{,}100~\textmu s & 20~\textmu s & $\approx\!14{,}900\times$ \\
\texttt{swap\_rows\_kernel} & 19 & 290{,}205~\textmu s & 25~\textmu s & $\approx\!11{,}600\times$ \\
\texttt{compute\_multipliers\_kernel} & 19 & 295{,}432~\textmu s & 20~\textmu s & $\approx\!14{,}800\times$ \\
\texttt{eliminate\_kernel} & 19 & 278{,}808~\textmu s & 20~\textmu s & $\approx\!13{,}900\times$ \\
\texttt{apply\_permutation\_kernel} & 1 & 300{,}320~\textmu s & --- & --- \\
\texttt{forward\_substitution\_kernel} & 1 & 326{,}736~\textmu s & --- & --- \\
\texttt{backward\_substitution\_kernel} & 1 & 324{,}597~\textmu s & --- & --- \\
\bottomrule
\end{tabular}
\caption{Per-kernel first-launch cost vs.\ typical (median) later-launch
cost, \texttt{TIME\_KERNELS=1} against \texttt{lu\_solve\_gpu} (8
kernels, well-conditioned \texttt{matrix.csv}). Every kernel's first
launch pays NVBit's one-time SASS-instrumentation-insertion cost
($\approx$280--330ms); every later launch of an already-instrumented
kernel costs $\approx$20--25 microseconds -- four to five orders of
magnitude faster.}
\label{tab:time-kernels-first-launch}
\end{table}

\emph{Every} kernel's first launch costs $\approx\!280$--$330$ms;
every later launch of an already-seen kernel costs $\approx\!20$--$25$
microseconds --- four to five orders of magnitude faster, a gap this
consistent across five structurally different kernels that it cannot be
explained by anything the kernels themselves do differently. The reason
is Section~\ref{sec:init}'s own finding, seen here as a direct timing
measurement rather than inferred from source: \texttt{instrument\_function}
rewrites a kernel's SASS \emph{only the first time that kernel function is
launched} (guarded by an \texttt{std::unordered\_set<CUfunction>} of
functions already instrumented), and \texttt{TIME\_KERNELS}'s clock
starts before that one-time rewrite and stops after it --- so the first
measurement for every kernel is dominated by SASS instrumentation-insertion
cost, not by anything the kernel computes. The three kernels launched only
once (\texttt{apply\_permutation\_kernel},
\texttt{forward\_substitution\_kernel}, \texttt{backward\_substitution\_kernel})
never get to show a fast, already-instrumented timing at all --- every
number \texttt{TIME\_KERNELS} ever reports for them is this one-time cost.
This is the same phenomenon at much larger scale behind the instrumentation
overhead \texttt{tc\_matmult/tc\_matmult\_report.tex} measures directly, and
behind this session's own attempt to run \texttt{pytorch\_slogdet/}
(elsewhere in this repository) under full instrumentation: cuSOLVER's
large, heavily-templated kernels made the one-time SASS rewrite alone take
long enough that a 120-second run was killed mid-instrumentation, before
the target program ever printed a line of its own output. Here, on a
20-row LU solve with only 8 small, simple kernels, the same cost is small
enough in absolute terms to see cleanly, kernel by kernel, in a single log
--- rather than dominating the whole run the way it does for a
library-heavy binary.

\subsection{\texttt{NIXNAN\_TIMEOUT}: Re-Enabled, Confirmed Active}
\label{sec:timeout-reenabled}

\textbf{Update:} an earlier version of this document, checked against an
earlier build on this same branch, found this timer disconnected and
called it dead. That finding is now stale. \texttt{git log -- src/timer.cu}
on the current checkout (\texttt{print\_histo\_instrn}) shows a commit
titled \texttt{Re-enable timer} sitting a handful of commits before the
tip --- the wiring below was reconnected after this document's original
``dead'' finding was written, and the rest of this section documents the
\emph{current}, active state, verified directly against a freshly rebuilt
\texttt{nixnan.so} rather than assumed from the earlier text.

\texttt{src/timer.h}/\texttt{timer.cu} implement a self-contained,
graceful internal timeout: \texttt{nixnan::timer::init()} reads
\texttt{NIXNAN\_TIMEOUT} (default 0, meaning disabled), and
\texttt{nixnan::timer::tool\_init(ctx)}, if the timeout is positive,
spawns a background thread that sleeps for that many seconds and then
calls \texttt{nnterminate("Timeout reached.")} --- a graceful,
in-process \texttt{SIGINT} to itself, flushing the log first.

\textbf{This code is called.} In \texttt{nixnan.cu}, the
init hook reads \texttt{NIXNAN\_TIMEOUT} via
\texttt{nixnan::timer::init()}, and the tool-init hook spawns the
timeout thread via \texttt{nixnan::timer::tool\_init(ctx)} --- both
confirmed directly in the current source, not inferred:
\begin{lstlisting}[language=C++]
void nvbit_at_init() {
  ...
  nixnan::fp_histogram::init();
  nixnan::timer::init();               // reads NIXNAN_TIMEOUT
}
void nvbit_tool_init(CUcontext ctx) {
  ...
  nixnan::timer::tool_init(ctx);       // spawns the timeout thread
}
\end{lstlisting}
This was confirmed two independent ways, both against a \texttt{nixnan.so}
rebuilt fresh from the current checkout:
\begin{itemize}
\item \textbf{The banner includes it.} \texttt{NIXNAN\_TIMEOUT} now
  appears in the ``NVBit core environment variables'' banner (between
  \texttt{BIN\_SPEC\_FILE} and \texttt{LOGFILE} in banner order --- see
  Table~\ref{tab:envvars}), which by the diagnostic established in
  Section~\ref{sec:init} means the running binary genuinely consults it.
\item \textbf{It actually terminates the process.} Run directly:
\begin{lstlisting}[language=bash]
NIXNAN_TIMEOUT=3 LD_PRELOAD=./nixnan.so ./some_long_running_target
\end{lstlisting}
produces \texttt{\#nixnan: Terminating early: Timeout reached.} in the
log at $\approx\!3.5$s, and the process exits with \texttt{SIGTERM}
(exit code 143) --- a real, graceful, in-process shutdown, not an
external kill.
\end{itemize}

The practical consequence: \texttt{bin/autonixnan}'s
\texttt{find\_extreme\_values()} sets \texttt{NIXNAN\_TIMEOUT} and
documents (in its own docstring) that this is what lets ``the target's
own timer, rather than the hard kill here'' end a long-running scan
gracefully --- and, on the current \texttt{nixnan.so}, it does exactly
that. The internal timer fires at \texttt{-t} seconds; the external
\texttt{subprocess.run(timeout=\ldots)} hard-kill is set to
\texttt{args.timeout + TIMEOUT\_GRACE} (30s) specifically so the graceful
internal exit has a 30-second window to complete --- flush its log,
tear down the GPU context --- before the external \texttt{SIGKILL} would
otherwise cut it off mid-shutdown. With the timer reconnected, that grace
period is no longer dead weight: it is now doing exactly the job its own
name describes.

\section{Part II --- \texttt{bin/autonixnan}}

\texttt{autonixnan} runs three phases against the same target command,
in sequence, unconditionally launching the target process fresh each
time (there is no persistent daemon or warm cache between phases).

\subsection{Phase 1: \texttt{collect\_kernels()} --- the baseline}

\begin{lstlisting}[language=Python]
env["LD_PRELOAD"] = lib_path
env["LOG_KERNELS"] = log_file
elapsed = run_with_env(program_args, timeout, env)
\end{lstlisting}

This runs the target with instrumentation fully disabled (Section
\ref{sec:logkernels}), so wall time here is close to native (modulo the
overhead of NVBit's own related-function enumeration and the
\texttt{"Function [\ldots] of kernel [\ldots]"} print spam for
library-heavy binaries, which this session measured at 10-17s for a
cuSOLVER-based binary that runs in under 2s uninstrumented). The
\texttt{-t}/\texttt{--timeout} value is used \emph{directly, with no
grace period}, as the hard \texttt{subprocess.run} timeout for this
phase --- a target that is merely slow to start (library init, matrix
setup) can therefore fail here even though nothing is actually being
instrumented, which this session hit directly with a cuSOLVER-heavy
binary whose factorization setup alone exceeded a naively small
\texttt{-t}.

The resulting log is parsed by \texttt{\_parse\_kernel\_log}, which only
matches lines starting \texttt{"\#nixnan: Kernel ["} --- ignoring the
\texttt{"Function [\ldots] of kernel [\ldots]"} noise described above ---
extracting the kernel name and, if present, its microsecond execution
time, and building both an in-order call sequence and a
name~$\to$~\{count, total\_time\} table.

\subsection{Report assembly (between phases 1 and 2)}

Before running Phase 2, \texttt{main()} prints the ``Unique Kernels''
table (sorted by total time descending, truncated to the top 25 when
there are more), the baseline wall-clock time, and the full kernel call
sequence --- elided to the first and last 25 calls with a
\texttt{"(\ldots calls omitted)"} marker in between when there are more
than 50 total, so an n=4096 factorization's 900+ launches do not flood
the terminal.

\subsection{Phase 2: \texttt{find\_extreme\_values()} --- the exponent scan}
\label{sec:phase2}

\begin{lstlisting}[language=Python]
bs = _make_binspec(1, max_reports)   # count=1: report every single hit, capped at max_reports/kernel
env["BIN_SPEC_FILE"] = spec_file
env["LOGFILE"] = log_file
env["NIXNAN_TIMEOUT"] = str(timeout)      # see Sec 2.9 -- confirmed active
env["INSTRUMENT_EXCEPTIONS"] = "0"        # exception checker off; histogram checker still runs
elapsed = run_with_env(program_args, timeout + TIMEOUT_GRACE, env, check=False)
\end{lstlisting}

\texttt{\_make\_binspec} auto-generates a \texttt{BIN\_SPEC\_FILE} that
watches the top 5\% and bottom 5\% of \emph{every} format's valid
(finite, non-zero) exponent range, for \emph{every} function --- since
the spec file has no per-function key, \texttt{fp-histogram.cu}'s
\texttt{process\_bin\_spec} applies the same global spec to whichever
function it is asked about, so in effect every instrumented kernel is
watched. This is a real, full NVBit exception-instrumentation-grade pass
over the entire binary (with the exception checker itself switched off
via \texttt{INSTRUMENT\_EXCEPTIONS=0}, leaving only the histogram
counter active) and is therefore subject to the same heavy per-distinct-kernel
overhead as any other full instrumentation run --- this session measured
it taking well over a minute for a cuSOLVER/cutlass-heavy binary with
hundreds of distinct kernel templates, versus milliseconds for a
handful of custom kernels called thousands of times.

\texttt{run\_with\_env} is called with \texttt{check=False}, so a
non-zero exit (which is the \emph{expected}, normal outcome here --- the
target hits \texttt{max\_reports} for some kernel and \texttt{nnterminate}s
itself with \texttt{SIGINT}) is not treated as an error; only an actual
\texttt{subprocess.TimeoutExpired} (after \texttt{timeout + 30}s) is
fatal to this phase.

The resulting log is scanned by \texttt{BIN\_REACHED\_RE} for
\texttt{"\#nixnan: <fmt> bin has reached threshold: function=\ldots
range=[lo,hi] count=N"} lines, tallied per function name, and the single
function with the most such reports across the whole run is chosen as
``the'' function of interest.

\subsection{Phase 3: \texttt{run\_instrumented()} --- the deep dive}

If Phase 2 found anything, \texttt{autonixnan} runs the target one final
time, with \texttt{LD\_PRELOAD} set and \emph{no other environment
variable} beyond a \texttt{FUNCTION\_WHITELIST} file containing exactly
the one function name Phase 2 flagged:

\begin{lstlisting}[language=Python]
env["LD_PRELOAD"] = lib_path
if whitelist_file:
    env["FUNCTION_WHITELIST"] = whitelist_file
elapsed = run_with_env(program_args, timeout, env, quiet=False)
\end{lstlisting}

Because no other variables are set, this run gets \texttt{INSTRUMENT\_EXCEPTIONS}'s
\emph{default} of 1 (full exception checking, Section~\ref{sec:exceptions})
and \texttt{HISTOGRAM}'s default of 0 (no histogram) --- but scoped, via
the whitelist, to \emph{only} the one function Phase 2 flagged (plus, per
\texttt{should\_instrument}, any function whose \emph{exact} name appears
in that whitelist; unrelated kernels are left completely uninstrumented,
running at native speed). \texttt{quiet=False} means this is the one
phase whose child process's \texttt{stdout}/\texttt{stderr} are not
suppressed --- this is why the live \texttt{\#nixnan: error [\ldots]}
lines and the final \texttt{------------ nixnan Report -----------}
aggregate (Section~\ref{sec:exceptions}) appear directly in
\texttt{autonixnan}'s own output during this phase and nowhere else.

This phase's own log file is opened and unconditionally
\texttt{os.unlink}ed in a \texttt{finally} block \emph{without ever being
read} --- \texttt{LOGFILE} is never set for this call (it is commented out
in the source), so the target's \texttt{nnout()} output goes to
\texttt{stderr}, which \texttt{quiet=False} already passes through to the
console directly; the temp file exists only so \texttt{tempfile.mkstemp}
has something to clean up, and \texttt{run\_instrumented}'s returned
\texttt{kernels\_seq}/\texttt{kernels\_unique} are consequently always
empty and are not used by \texttt{main()}.

\textbf{This is also the phase where the exponent-scan/exception-checker
scope mismatch documented in Section~\ref{sec:histogram} bites hardest}:
if the real fault is a hard jump straight to \texttt{Inf} (rather than a
gradual climb through extreme-but-finite exponents), Phase 2 will not
flag the kernel where it actually happens, so Phase 3's whitelist will
scope full exception-instrumentation to the \emph{wrong} kernel, and the
final aggregate report can show all-zero counts in every format even
though the target program's own internal exception detection (if it has
any) clearly fired. Section~\ref{sec:limitations} gives the exact
reproduction from this session.

\subsection{Timeout semantics, summarized}

\begin{center}
\begin{tabular}{lll}
\toprule
Phase & Hard kill after & Notes \\
\midrule
1: baseline & \texttt{-t} exactly & No grace period; catches slow startup too. \\
2: exponent scan & \texttt{-t + 30s} & Grace period for the internal timer to shut down gracefully first (Sec.~\ref{sec:timeout-reenabled}). \\
3: deep dive & \texttt{-t} exactly & Same value re-used; not separately configurable. \\
\bottomrule
\end{tabular}
\end{center}

Every phase re-uses the \emph{same} \texttt{-t} value from the command
line; there is no way to give the baseline phase a short budget and the
scan phase a long one from the CLI as it stands.

\section{Part III --- \texttt{iter-autonixnan} and \texttt{ksel-autonixnan}}
\label{sec:iter-ksel}

Both tools were added to \texttt{bin/} to sidestep Phase 2's
exponent-scan blind spot (Section~\ref{sec:phase2},
Section~\ref{sec:limitations} item 1) entirely, by never depending on it
to pick a kernel in the first place. Neither reads or writes a
\texttt{BIN\_SPEC\_FILE} unless the caller supplies one with \texttt{-b};
by default both simply run every kernel they touch under full exception
checking (\texttt{INSTRUMENT\_EXCEPTIONS=1}, forced, regardless of the
ambient environment).

\subsection{\texttt{iter-autonixnan}: one kernel at a time}

\begin{lstlisting}[language=bash]
iter-autonixnan -s SAMPLING [-t TIMEOUT] [-b BIN_SPEC_FILE] -- PROGRAM [ARGS...]
\end{lstlisting}

Phase 1 is the same \texttt{collect\_kernels()} baseline as
\texttt{autonixnan}'s own Phase 1 (Section~\ref{sec:phase2}'s sibling,
re-implemented rather than imported). Phase 2 then loops over every
kernel name \texttt{k} found, and for each one, in turn: writes a
\texttt{FUNCTION\_WHITELIST} file containing \emph{only} \texttt{k}, and
launches an entirely fresh process with \texttt{SAMPLING=<-s value>},
\texttt{INSTRUMENT\_EXCEPTIONS=1}, and \texttt{BIN\_SPEC\_FILE=<-b
value>} if one was given. Every other kernel is completely uninstrumented
for that run (native speed); the target program runs to completion (or
its own timeout) once per kernel in \texttt{KN}.

This gives clean, unambiguous per-kernel attribution: each kernel's
exceptions are reported in complete isolation from every other kernel's.
The cost is $|KN| + 1$ full process launches of the target program.
For a target with a handful of custom kernels (e.g.\
\texttt{lu\_solve\_gpu}, 8 kernels; \texttt{rd\_nixnan}, 2--4 kernels)
this is cheap. For a target whose kernels come from a large, heavily
templated library (PyTorch, cuSOLVER, cuBLAS), it can be
\textbf{catastrophically expensive} -- see
Section~\ref{sec:limitations} item 4, where sweeping just 14 kernels
this way was projected at roughly 6.5 hours and abandoned partway
through.

\subsection{\texttt{ksel-autonixnan}: a hand-picked combination, together}

\begin{lstlisting}[language=bash]
ksel-autonixnan -phase1 [--ksel-file FILE] [-t TIMEOUT] -- PROGRAM [ARGS...]
ksel-autonixnan -ksel [--ksel-file FILE] -s SAMPLING [-t TIMEOUT] [-b BIN_SPEC_FILE] -- PROGRAM [ARGS...]
\end{lstlisting}

\texttt{-phase1} and \texttt{-ksel} are mutually exclusive and one is
required. \texttt{-phase1} runs the same kernel-collection baseline as
\texttt{iter-autonixnan}'s own Phase 1, then writes a \emph{selection
file} (default path \texttt{./ksel}) -- one line per kernel, of the form

\begin{lstlisting}
<exact kernel name> : n ;
\end{lstlisting}

every kernel defaulted to \texttt{n} (not selected). The intended
workflow is to hand-edit this file, flipping \texttt{n} to \texttt{y} for
whichever kernels should be examined, then re-run with \texttt{-ksel}.
\texttt{parse\_ksel\_file} splits each line on the \emph{last}
\texttt{" : "} occurrence (via a regex that requires no further
\texttt{:} after the match, \texttt{\textbackslash s:\textbackslash
s(?=[\^{}:]*\$)}), specifically so that a kernel name itself containing
C++'s \texttt{::} scope operator is never mistaken for the field
separator.

\texttt{-ksel} collects every kernel marked \texttt{y}, writes them
\emph{all into one} \texttt{FUNCTION\_WHITELIST} file, and launches the
target \textbf{exactly once}, with all of them whitelisted together for
that single run. Unlike \texttt{iter-autonixnan}, there is no
per-kernel isolation: if two selected kernels both produce exceptions,
the final aggregate report pools them together indistinguishably.

The payoff is cost: whatever one-time cost NVBit's fatbinary enumeration
carries is paid \textbf{once}, for the single combined process, no
matter how many kernels are selected. Section~\ref{sec:limitations}
item 4 gives the direct head-to-head measurement -- for the same 14
kernels that were projected to need $\sim$6.5 hours under
\texttt{iter-autonixnan}, the \texttt{ksel-autonixnan} combined run
completed in 1720.1 seconds (about the same as \emph{one} kernel's own
enumeration cost under \texttt{iter-autonixnan}), and found real
exceptions in the process.

\subsection{Choosing between the three}

Table~\ref{tab:choosing-three} summarizes the tradeoffs.

\begin{table}[p]
\centering
\small
\begin{tabular}{p{2.1cm}p{3.0cm}p{2.6cm}p{2.6cm}p{3.4cm}}
\toprule
& \textbf{\texttt{autonixnan}} & \textbf{\texttt{iter-autonixnan}} & \textbf{\texttt{ksel-autonixnan}} & \textbf{direct \texttt{nixnan.so}} \\
\midrule
\textbf{Kernel selection} &
  automatic (Phase 2's exponent scan picks one) &
  every kernel, one at a time &
  hand-picked subset, all together &
  everything (or a manual whitelist/blacklist) \\
\textbf{Process launches} &
  3 (baseline, scan, deep-dive) &
  $|KN|+1$ &
  2 (\texttt{-phase1} then \texttt{-ksel}) &
  1 \\
\textbf{Exception attribution} &
  single kernel (whichever Phase 2 flagged) &
  per-kernel, fully isolated &
  pooled across the selected subset &
  pooled across everything instrumented \\
\textbf{Dominant cost} &
  Phase 2's full histogram-scan pass over the whole binary &
  fatbinary-enumeration cost paid $|KN|+1$ times &
  fatbinary-enumeration cost paid once &
  fatbinary-enumeration cost paid once \\
\textbf{Blind spot} &
  hard finite-to-Inf/NaN jumps invisible to Phase 2 (item 1 below); wrong kernel can get whitelisted &
  none from exponent-scan bias; only the enumeration-cost blowup (item 4) &
  same as \texttt{iter-autonixnan}, plus needing a human to pick the subset (or select everything, degrading to direct) &
  none from selection bias; must fit the full instrumentation cost in one run \\
\textbf{Best for} &
  a first, unattended look at an unfamiliar binary &
  a handful of small/custom kernels needing individual attribution &
  a library-heavy binary where enumeration cost dominates and pooled attribution is acceptable &
  anything small/fast enough to instrument whole, with no scan-bias risk \\
\bottomrule
\end{tabular}
\caption{Choosing between \texttt{autonixnan}, \texttt{iter-autonixnan},
\texttt{ksel-autonixnan}, and a direct, unwrapped \texttt{nixnan.so}
run -- kernel selection strategy, process-launch count, exception
attribution granularity, dominant cost, blind spot, and best use case
for each.}
\label{tab:choosing-three}
\end{table}

\section{Empirically Observed Limitations}
\label{sec:limitations}

These are not code-reading conclusions but things directly reproduced
against real binaries this session, cross-checked against the source
above.

\begin{enumerate}
\item \textbf{Hard overflow jumps are invisible to the extreme-value
  scan.} Running a mixed-precision Cholesky iterative-refinement solver
  (\texttt{SC26-expts/mpir\_gpu}) with its documented overflow-forcing
  flags (\texttt{--noscale --bscale 1e5}) produced a program-confirmed
  FP16 overflow (\texttt{nonfinite=256} from the target's own internal
  check) on the very first iteration. \texttt{autonixnan}'s Phase 2 never
  flagged the kernel responsible; it flagged an unrelated fp64 QR kernel
  used only during test-matrix generation, and Phase 3's whitelisted
  full-instrumentation pass consequently reported zero FP16 exceptions.
  Running \texttt{nixnan.so} directly (no \texttt{autonixnan}, no
  whitelist, default \texttt{INSTRUMENT\_EXCEPTIONS=1}) against the same
  invocation found the real exceptions immediately --- filed under the
  \texttt{FP32} bucket, not \texttt{FP16}, because that solver
  configuration only rounds values into fp16's representable range via a
  cast; the arithmetic that actually overflows executes as native FP32
  SASS instructions. Neither of these is a nixnan bug: the histogram
  mechanism's zero/inf exclusion is by design (Section~\ref{sec:histogram}),
  and the FP32 attribution is simply an accurate report of which SASS
  instruction type executed on the bad value.
\item \textbf{The internal timeout, at the time of the runs this
  session's numbers come from}, was disconnected --- every \texttt{-t}
  used here was enforced purely externally, with a (then-)pointless 30s
  grace period tacked onto the scan phase. \textbf{This is no longer
  true}: a later check against a freshly rebuilt \texttt{nixnan.so} on
  this same branch found the timer reconnected (a
  \texttt{Re-enable timer} commit sits between the two checks) and
  genuinely terminating the process on schedule --- see
  Section~\ref{sec:timeout-reenabled} for the direct verification. This
  item is kept as an accurate record of what this session's own runs
  actually exercised, not as a current statement about the tool.
\item \textbf{\texttt{LOG\_KERNELS} noise scales with library surface
  area, not with actual kernel diversity.} A cuSOLVER/cuBLAS-heavy binary
  produced 10{,}000+ log lines from related-function enumeration alone,
  none of which affect \texttt{autonixnan}'s parsed output but which do
  make the raw log file unwieldy to read by hand.
\item \textbf{NVBit's one-time fatbinary-enumeration cost is paid per
  \emph{process}, not per \emph{instrumented kernel}.}
  (\texttt{Karpathy-NGram/}, full data in that repository's
  \texttt{All-Runs.md}.) \texttt{gpt\_finite\_\allowbreak state\_cuda.py} -- a trivial,
  12{,}656-parameter toy GPT, reduced to a single forward pass -- was run
  under \texttt{iter-autonixnan}. Phase 1 found 14 distinct kernels this
  one forward pass touches and took 1707.7s just to enumerate them once,
  uninstrumented. Phase 2's very first per-kernel whitelisted run (the
  \texttt{arange} index kernel alone) took \emph{another} 1713.9s on its
  own -- essentially the same cost all over again, for a single kernel,
  with 13 more still to go. At that rate the full 14-kernel sweep was on
  pace for roughly 6.5 hours and was abandoned. Switching to
  \texttt{ksel-autonixnan} (Section~\ref{sec:iter-ksel}) -- whitelisting
  all 14 kernels together in \emph{one} combined process instead of 14
  separate ones -- completed the equivalent work in 1720.1s: essentially
  identical to the cost of enumerating \emph{one} kernel under
  \texttt{iter-autonixnan}, roughly a 13x wall-time reduction for the
  same 14-kernel coverage, and it found real exceptions doing so (71
  distinct FP32 NaN sites, 22 distinct FP32 -Infinity sites, 6 distinct
  FP16 NaN sites, inside PyTorch's own
  \texttt{softmax\_warp\_\allowbreak forward} kernel -- traced to the
  causal attention mask's \texttt{-inf} values
  interacting with the kernel's \texttt{exp2}-based, Kahan-split
  \texttt{exp(x)} decomposition; the model's own final output was
  unaffected, since softmax of \texttt{-inf} legitimately evaluates to
  \texttt{0}). This confirms the dominant cost is NVBit's one-time
  analysis of the \emph{entire loaded CUDA fatbinary} per fresh process,
  independent of how many (or how few) kernels are actually whitelisted
  for instrumentation -- a real, practical reason to prefer
  \texttt{ksel-autonixnan}'s combined-whitelist strategy over
  \texttt{iter-autonixnan}'s per-kernel sweep for any library-heavy
  binary (PyTorch, cuSOLVER, cuBLAS) needing more than one or two
  kernels examined.
\item \textbf{NVBit-based instrumentation is not fully robust at the
  scale and architectural complexity of a real, production-sized LLM.}
  (\texttt{biomistral-test/}, full report and logs in that directory.)
  Attaching \texttt{nixnan.so} (\texttt{SAMPLING=512}) to a
  \textbf{7-billion-parameter}, 4-bit-quantized (\texttt{bitsandbytes},
  NF4) medical dialog model (BioMistral-7B) for a single conversational
  turn reproducibly crashed with \texttt{SIGSEGV} before ever reaching
  \texttt{nixnan}'s own final aggregate report, in every one of 4
  attempts. Two of those attempts were the \emph{identical} command run
  twice, to check for a transient race: they reproduced almost exactly
  (45 distinct kernels touched both times, near-identical exception
  counts, and the identical last-logged kernel before the crash),
  confirming this is deterministic, not a fluke. \texttt{SAMPLING} does
  not mitigate it, since sampling only gates the runtime per-launch
  enable/disable toggle (Section~\ref{sec:sampling-launch-gate}), not the
  one-time SASS rewrite where the fault appears to occur. Progressively
  adding the last-logged kernel before each crash to a
  \texttt{FUNCTION\_BLACKLIST} pushed the crash substantially later each
  round -- from 45 distinct kernels successfully touched, to 91, reaching
  well into \texttt{generate()}'s own token-sampling logic
  (\texttt{ArgMaxOps}, an exponential-distribution sampling kernel) --
  without ever reaching completion after 4 rounds. A proper stack trace
  was attempted via Ubuntu's \texttt{apport} crash reporting, but the
  resulting core dump was 57\,GB (a 7B-parameter model's GPU-mapped
  memory counts toward the crashing process's own address space),
  impractical to load into \texttt{gdb} in this environment. Despite
  never completing, every attempt captured substantial, genuine exception
  data (500+ distinct NaN sites, roughly 300 distinct -Infinity sites,
  roughly 230--300 distinct subnormal sites per run, across $\sim$90
  kernels once the two known crash-adjacent kernels were blacklisted) --
  this is reported as a real reliability boundary of the tool at this
  scale, not a defect specific to this model or input, and stands in
  contrast to every other (much smaller-scale) case study in this
  repository completing cleanly under the same \texttt{nixnan.so} build.
\end{enumerate}

\section{Practical Summary}

To reproduce any given phase by hand outside \texttt{autonixnan}, the
minimal commands are:

\begin{lstlisting}[language=bash]
# Phase 1 equivalent -- kernel catalog, no instrumentation
LD_PRELOAD=./nixnan.so LOG_KERNELS=/tmp/k.log ./target [args...]

# Phase 2 equivalent -- exponent-extreme scan, exceptions off
LD_PRELOAD=./nixnan.so BIN_SPEC_FILE=/tmp/spec.json LOGFILE=/tmp/scan.log \
  INSTRUMENT_EXCEPTIONS=0 ./target [args...]

# Phase 3 equivalent, or a direct full exception run (no whitelist
# needed if the binary is small/fast enough to instrument in full)
LD_PRELOAD=./nixnan.so ./target [args...]
\end{lstlisting}

The last form -- direct, unwhitelisted \texttt{nixnan.so}, no
\texttt{autonixnan} involved -- is the one to reach for whenever a target
is small enough to instrument in full and the exponent-scan's blind spot
(hard overflow jumps) is a concern, since it is the only one of the three
that does not depend on Phase 2 having correctly identified the kernel of
interest.

When that direct run is impractical (too many kernels to instrument all
at once, or the process needs to survive with only a subset instrumented),
reach for \texttt{iter-autonixnan} or \texttt{ksel-autonixnan}
(Section~\ref{sec:iter-ksel}) instead of \texttt{autonixnan} itself --
both sidestep Phase 2's exponent-scan blind spot entirely, by never
depending on it to pick a kernel:

\begin{lstlisting}[language=bash]
# Every kernel, one at a time, fully isolated attribution -- cheap for a
# handful of small/custom kernels, ruinously expensive for a
# library-heavy binary (Section~\ref{sec:limitations} item 4)
iter-autonixnan -s SAMPLING -- ./target [args...]

# A hand-picked subset, all together in one combined run -- the
# enumeration cost is paid once regardless of subset size
ksel-autonixnan -phase1 -- ./target [args...]   # writes ./ksel; hand-edit it
ksel-autonixnan -ksel -s SAMPLING -- ./target [args...]
\end{lstlisting}

\end{document}
